\label{sec:related_work}

\subsection{TBM response modelling and spatial gaps}
Traditional TBM performance studies estimate penetration, advance rate,
thrust, torque, or related operational responses from rock mass properties and
machine parameters. Empirical, statistical, and machine-learning formulations
are used for planning and performance assessment in hard-rock tunnelling
\citep{yagiz2008utilizing,hassanpour2010tbm,armaghani2018prediction}. With the
accumulation of monitoring data, sequence models such as LSTM and GRU also
learn response inertia and lagged dependence from recent operational histories
\citep{hochreiter1997long,cho2014learning}. These studies provide important
response-prediction baselines, but their spatial unit is usually a tunnel
section, a sample, or a feature vector rather than an explicit relation between
a geological object and a TBM component.

This limitation is central to the present paper. A monitoring curve can
describe how torque, thrust, penetration, advance rate, or shield pressure
changes with time or chainage, but it does not state which geological voxels
are plausible interaction partners for the cutterhead, front shield, middle
shield, or tail shield. The present study therefore does not try to replace
monitoring-based predictors. It asks whether a spatial relation representation
provides component-indexed field information about the rock--machine
configuration that monitoring curves alone do not encode.

\subsection{Geological prediction as spatial support}
Advance geological prediction methods, including TSP and related geophysical
or drilling-based approaches, are used to identify faults, fractured zones,
water-rich zones, weak surrounding rock, and other geological anomalies before
or during tunnel excavation \citep{li2019tunnel}. For TBM analysis, such data
are valuable because they describe spatial heterogeneity that is not contained
in the monitoring record itself. TSP results can be represented as velocity
profiles, sections, or voxelised attribute fields along the tunnel direction;
voxelisation is useful here because it preserves spatial support, local
attribute values, and chainage-referenced geological context.

Geological spatial representation alone is nevertheless incomplete for the
diagnostic question addressed in this paper. A voxel field can show where
anomalous ground occurs, but it does not specify whether that anomaly is
adjacent to the cutterhead or shield, whether the adjacency is
normal-compatible, or whether the voxel remains active at the current
excavation step. This is a spatial-support problem as much as an engineering
problem: TSP voxels, TBM surface samples, and monitoring records have
different spatial or observational supports and cannot be interpreted as a
single tabular unit without losing relation information
\citep{goodchild1992geographical,longley2005geographic,openshaw1984modifiable}.

\subsection{Entity--relation spatial data models}
GIScience is concerned not only with mapping locations but also with
representing geographical objects, spatial relations, temporal change, and
processes \citep{goodchild1992geographical,worboys2004gis,galton2001spatial}.
Spatial relation theory provides concepts for neighbourhood, adjacency,
containment, intersection, and other topological or metric relations
\citep{egenhofer1991point,wu2008contiguity}. TBM excavation fits this
perspective because rock units, machine surfaces, and monitoring responses are
observed in a moving excavation coordinate, and the relevant relations change
as the machine advances.

Recent geospatial knowledge-graph and graph data-model studies show how
heterogeneous spatial entities are integrated for retrieval, reasoning, and
query \citep{zhu2026knowwheregraph,xu2025spatiotemporal,li2025question,rahimi2021topology}.
Scenario-oriented spatiotemporal knowledge-graph models likewise emphasise
linking heterogeneous data through meaningful objects and relations rather
than storing them only as parallel indicator tables
\citep{wang2025urbanphysical}. The present paper shares this entity--relation
modelling perspective, but it operates at a finer voxel and surface-node
support and targets a moving underground engineering process rather than an
indoor social network, an urban knowledge service, or a regional knowledge
base.

This literature provides the conceptual foundation for the proposed graph:
rock voxels are geological objects, TBM surface samples are machine-side
spatial objects, and rock--machine edges encode screened candidate interaction
relations. Monitoring responses are kept as aligned records rather than graph
nodes. The representational gap is therefore not the absence of a predictor,
but the lack of a spatial data model that connects geological voxels, machine
components, candidate interaction relations, and response residuals in one
component-indexed field representation.

\subsection{Spatial interaction graphs}
Graphs are a natural formalism for representing irregular spatial objects and
their relations. Recent geospatial studies have used graph-based
representations for spatial classification, spatial multivariate
distributions, urban network morphology, and spatially constrained community
detection
\citep{zhu2020spatial,desabbata2023graph,liu2025urban,chen2026usgt}. Spatial
interaction models further emphasise that relations are shaped by distance
decay, compatibility, direction, accessibility, and spatial support
\citep{wilson1971spatial,fotheringham1989spatial}. For the rock--TBM problem,
this means that candidate relations should not be fully connected or purely
learned without physical constraints. They should depend on active-zone
membership, distance, surface-normal compatibility, and excavation state.

Graph neural networks and heterogeneous graph models remain relevant as future
extensions for learning on structured data
\citep{schlichtkrull2018modeling,zhang2019heterogeneous,wu2020comprehensive,zhou2020graph}.
Spatiotemporal graph models have also been successful in traffic, flow
prediction, and physical simulation
\citep{li2018diffusion,wu2019graph,zhang2020deep,geng2019spatiotemporal,sanchez2020learning,pfaff2021learning}.
However, these studies mainly motivate the representational capacity of graph
structures and the need for spatially constrained relations. They do not imply
that the present tunnel cases should be framed primarily as a black-box
graph-learning benchmark.

In the tunnelling domain, digital twin frameworks for TBM performance
prediction, visualisation, monitoring, and deformation prediction have been
proposed \citep{latif2023digital,apoji2023shaping,hu2025dataknowledge}. The
present graph-based representation is complementary: rather than replacing
prediction-oriented digital twin modules, it specifies an underground spatial
relation layer between geological voxels and TBM surface components. This layer
defines active candidate relations, component-indexed aggregation, and
edge-level traceability, which are typically implicit in monitoring-centred
digital twin workflows. The evaluation is therefore kept descriptive:
Spearman association, detrending, null comparisons, sensitivity analysis, and
edge-level traceability are used to support a spatial field-representation
claim, not a universal prediction or risk-calibration claim.
